\chapter{Memory Hierarchy and Caches}\label{ch:ll:memory-hierarchy-and-caches}

On this book's laptop, reading one price level of an order book costs about a nanosecond when the level is in the
core's first-level cache, three or four nanoseconds when it is in the second level, and about 150 nanoseconds when it
has to come from memory: a factor of more than a hundred for the same instruction. Chapter~1's toy book, built on balanced trees
whose nodes are scattered across the heap, spent most of its time waiting for such loads. Most of the difference
between a fast order book and a slow one is not the algorithm but where its bytes live, and this chapter is about that:
the levels of the hierarchy and their costs, what happens when two cores share a line, how virtual addresses are
translated, and what the hardware does to hide all of it.

\section{Cache levels, lines and associativity}

\begin{definition}[Cache line, set-associative cache]\label{def:ll:memory-hierarchy-and-caches:line}
Memory moves between the levels of the hierarchy in fixed blocks called \emph{cache lines}\index{cache line}, 64 bytes
on current x86 processors. In a \emph{set-associative cache}\index{set-associative cache} with $S$ sets of $W$ ways, a
line can be stored only in the set given by bits of its address, in any of that set's $W$ slots; the cache holds
$S \cdot W$ lines.
\end{definition}

\begin{definition}[Cache miss, working set]\label{def:ll:memory-hierarchy-and-caches:miss}
A \emph{cache miss}\index{cache miss} is an access to a line that is not in a given level, which must then be fetched
from the next level down. The \emph{working set}\index{working set} of a piece of code over an interval is the set of
lines it touches in that interval; the level whose capacity it fits in sets its speed.
\end{definition}

An address splits into three fields (\cref{fig:ll:memory-hierarchy-and-caches:split}): the low 6 bits are the offset
within the 64-byte line; the next $\log_2 S$ bits choose the set; the rest, the tag, identify the line among those that
map to the set. Two consequences for data layout. A structure that straddles two lines costs two misses where one would
do. And addresses that differ by a multiple of $S \times 64$ bytes compete for the same $W$ slots: a program that walks
several arrays whose strides are that multiple (a power of two, typically 4\,KiB) can miss in a cache that is mostly
empty.

\begin{omfigure}
\begin{tikzpicture}[font=\footnotesize]
\draw[fill=gray!10] (0,0) rectangle (6.4,0.6); \node at (3.2,0.3) {tag};
\draw[fill=omDef!15] (6.4,0) rectangle (8.8,0.6); \node at (7.6,0.3) {set index};
\draw[fill=omThm!12] (8.8,0) rectangle (10.8,0.6); \node at (9.8,0.3) {offset};
\node[anchor=south,align=center] at (3.2,0.65) {remaining bits of the address};
\node[anchor=south,align=center] at (7.6,0.65) {bits 6--11\\ (64 sets)};
\node[anchor=south,align=center] at (9.8,0.65) {bits 0--5\\ (64 bytes)};
\draw[-Stealth,thick,omDef] (7.6,0) -- (7.6,-0.75);
\foreach \w in {0,...,11} \draw[draw=omDef,fill=omDef!6] (4.6+0.5*\w,-0.8) rectangle (5.05+0.5*\w,-1.3);
\node[anchor=north,align=center] at (7.35,-1.4) {the selected set: 12 ways, any of which may hold a line with this index};
\end{tikzpicture}

\omcaption{How an address selects its place in a \omterm{def:ll:memory-hierarchy-and-caches:line}{set-associative cache}. A 48~KiB, 12-way cache of 64-byte lines has 64
sets; bits 6 to 11 of the address choose the set, so addresses 4\,KiB apart share a set.}\label{fig:ll:memory-hierarchy-and-caches:split}
\end{omfigure}

\begin{example}[Geometry of a first-level cache]\label{ex:ll:memory-hierarchy-and-caches:geometry}
A 48~KiB, 12-way cache with 64-byte lines has $48 \times 1024 / (64 \times 12) = 64$ sets. The set index is bits 6 to
11; addresses that are equal modulo $64 \times 64 = 4\,096$ bytes fall into the same set. Thirteen hot structures placed
at the same offset of thirteen 4~KiB pages cannot all stay in this cache, although together they occupy 832 bytes.
\end{example}

\Cref{fig:ll:memory-hierarchy-and-caches:stairs} measures the hierarchy with a pointer chase: a buffer of $n$ slots, one
per \omterm{def:ll:memory-hierarchy-and-caches:line}{cache line}, laid out as a single random cycle, each slot holding the address of the next. Each load's address is
the previous load's result, so no two loads overlap and the time per step is the latency of the level the buffer fits
in. The staircase shows the first level up to 48~KiB at about \qty{1.1}{\nano\second} (five cycles), the second
level up to about a megabyte and a half at about \qty{3.5}{\nano\second}, the third level from about three to eight
megabytes at 15 to \qty{23}{\nano\second}, close to its published latency of 75 cycles or more (about
\qty{20}{\nano\second}), and memory from about twelve megabytes on, at 100 to \qty{175}{\nano\second}: the
low-power memory of a laptop, which a published measurement puts at 150 to 175 nanoseconds for this processor. The
third level is 24~MiB, shared by all the cores; a random chase through a virtual machine's pages starts to miss in it
well before its nominal size.

\begin{omfigure}
\begin{tikzpicture}
\begin{axis}[width=12cm,height=6cm,xmode=log,ymode=log,log basis x=2,xmin=4096,xmax=268435456,ymin=0.8,ymax=500,
  xtick={4096,32768,262144,2097152,16777216,134217728},xticklabels={4 KiB,32 KiB,256 KiB,2 MiB,16 MiB,128 MiB},
  ytick={1,2,5,10,20,50,100,200},yticklabels={1,2,5,10,20,50,100,200},
  xlabel={working set (log scale)},ylabel={ns per load (log scale)},
  tick label style={font=\small},label style={font=\small},grid=major,grid style={gray!25},
  legend columns=3,legend style={font=\footnotesize,at={(0.5,-0.3)},anchor=north,draw=none,fill=none,
  /tikz/every even column/.append style={column sep=8pt}},table/col sep=comma]
\addplot[omDef,very thick,mark=*,mark size=1.2pt] table[x=bytes,y=random_ns]{figdata/low-latency/03-memory-hierarchy-and-caches/measured_stairs.csv};
\addplot[omThm,thick,mark=square*,mark size=1.2pt] table[x=bytes,y=random_huge_ns]{figdata/low-latency/03-memory-hierarchy-and-caches/measured_stairs.csv};
\addplot[omMeth,thick,dashed,mark=triangle*,mark size=1.2pt] table[x=bytes,y=sequential_ns]{figdata/low-latency/03-memory-hierarchy-and-caches/measured_stairs.csv};
\legend{random order,{random order, huge pages},address order}
\end{axis}
\end{tikzpicture}

\omcaption{The memory staircase: time per dependent load through a buffer of one slot per \omterm{def:ll:memory-hierarchy-and-caches:line}{cache line}, against the
buffer's size. In address order the prefetcher hides almost all of it. Measured on a laptop (Intel Core Ultra 7 155H)
under WSL2, no isolated cores, pinned to one CPU, the machine otherwise idle. Data: \texttt{bench\_mem.py}.}\label{fig:ll:memory-hierarchy-and-caches:stairs}
\end{omfigure}

\section{Coherence and false sharing}

\begin{definition}[Cache coherence]\label{def:ll:memory-hierarchy-and-caches:coherence}
\emph{Cache coherence}\index{cache coherence} is the protocol by which the private caches of several cores keep one
view of each line: before a core writes a line it must obtain the line in an exclusive state, invalidating every other
core's copy; a later read by another core must fetch the modified line from the writer. In the MESI family each line
in each cache is Modified, Exclusive, Shared or Invalid.
\end{definition}

\begin{definition}[False sharing]\label{def:ll:memory-hierarchy-and-caches:false}
\emph{False sharing}\index{false sharing} occurs when two cores repeatedly write different variables that lie in the
same \omterm{def:ll:memory-hierarchy-and-caches:line}{cache line}: neither shares data with the other, yet every write invalidates the other core's copy and the line
travels between them.
\end{definition}

A market-data thread that updates a sequence counter and a strategy thread that updates its own statistics, placed next
to each other in one structure, pay the price of the line moving between cores on every update, without any data
being shared. \Cref{fig:ll:memory-hierarchy-and-caches:false} measures it: one thread incrementing a counter pays about
\qty{4}{\nano\second} per atomic increment; two threads each incrementing their own counter on separate lines pay about
the same; two threads whose counters share a line pay about five times as much. The cure is layout: give each
independently written variable its own line (\texttt{CachePadded}, below), and keep variables written by one thread
together and away from those written by another.

\begin{omfigure}
\begin{tikzpicture}
\begin{axis}[width=9cm,height=5cm,xbar,bar width=12pt,xmin=0,enlarge x limits={upper,value=0.15},
  symbolic y coords={same line,padded,one thread},ytick=data,enlarge y limits=0.3,
  xlabel={ns per atomic increment},tick label style={font=\small},label style={font=\small},
  grid=major,grid style={gray!25},nodes near coords,nodes near coords style={font=\footnotesize},
  every node near coord/.append style={/pgf/number format/fixed,/pgf/number format/precision=1},
  table/col sep=comma]
\addplot[fill=omThm!45,draw=omThm] table[y=layout,x=ns_per_increment]{figdata/low-latency/03-memory-hierarchy-and-caches/measured_false.csv};
\end{axis}
\end{tikzpicture}

\omcaption{\omterm{def:ll:memory-hierarchy-and-caches:false}{False sharing}: two threads, pinned to two CPUs, each incrementing its own atomic counter 20 million times,
with the counters on one \omterm{def:ll:memory-hierarchy-and-caches:line}{cache line} or on two. Measured on a laptop (Intel Core Ultra 7 155H) under WSL2, no isolated
cores. Data: \texttt{bench\_mem.py}.}\label{fig:ll:memory-hierarchy-and-caches:false}
\end{omfigure}

\section{Translation buffers and huge pages}

\begin{definition}[Translation lookaside buffer, huge page]\label{def:ll:memory-hierarchy-and-caches:tlb}
Programs use virtual addresses, translated to physical addresses through page tables. The \emph{translation lookaside
buffer}\index{translation lookaside buffer} (TLB) caches recent translations; a TLB miss makes the processor walk the
page tables, several dependent memory accesses. A \emph{huge page}\index{huge page} is a page larger than the base
4~KiB (2~MiB or 1~GiB on x86-64), which one TLB entry covers.
\end{definition}

The TLB's reach is its number of entries times the page size. A second-level TLB of 2\,048 entries covers 8~MiB of
4~KiB pages and 4~GiB of 2~MiB pages; an order book, a symbol table and a few rings of a trading process can easily
exceed the first and never the second. Linux offers \omterm{def:ll:memory-hierarchy-and-caches:tlb}{huge pages} two ways: reserved pages from a pool, and transparent
\omterm{def:ll:memory-hierarchy-and-caches:tlb}{huge pages}, which the kernel assembles on its own; in the default \texttt{madvise} mode it does so only for regions the
program has marked with \texttt{madvise(MADV\_HUGEPAGE)}. The kernel's documentation notes that the gain is larger
under virtualisation, where each walk goes through two levels of page tables, the guest's and the hypervisor's. On
this laptop, which is virtualised, a random chase over 256~MiB costs about \qty{180}{\nano\second} a step with 4~KiB
pages and about \qty{150}{\nano\second} with \omterm{def:ll:memory-hierarchy-and-caches:tlb}{huge pages} (\texttt{measured\_huge.csv}): the difference is
the page walk.


\section{Prefetching and access patterns}

\begin{definition}[Hardware prefetching]\label{def:ll:memory-hierarchy-and-caches:prefetch}
\emph{Hardware prefetching}\index{hardware prefetching} is the processor's fetching of lines before they are requested,
by recognising patterns in recent misses: consecutive lines, constant strides.
\end{definition}

The dashed curve of \cref{fig:ll:memory-hierarchy-and-caches:stairs} is the same chase with the slots visited in address
order: the prefetcher recognises the stream and the loads cost under \qty{2}{\nano\second} until the buffer leaves the
second level, and four or five nanoseconds from memory, some thirty times less than the random chase. Nothing changed but
the order. The practical rules follow: store what the hot path reads together in contiguous arrays, not in nodes linked
by pointers; iterate in address order; keep the hot data of a message (price, quantity, side) in the first line of its
structure; and when the access is truly random, as an order-id lookup is, make the table small enough to stay in cache
(chapter~19).

\section{Tutorial: measuring the staircase}\label{tut:ll:memory-hierarchy-and-caches:tut}

\begin{tutorial}
\textbf{Goal.} Measure the latency of each level of the hierarchy, the cost of \omterm{def:ll:memory-hierarchy-and-caches:false}{false sharing} and the effect of \omterm{def:ll:memory-hierarchy-and-caches:tlb}{huge
pages}. \textbf{End state:} \cref{fig:ll:memory-hierarchy-and-caches:stairs,fig:ll:memory-hierarchy-and-caches:false} and
the huge-page table.
\begin{tutsteps}
\item \textbf{A ring of pointers.} Shuffle the slot order with Fisher--Yates, then store in each slot the address of the
next: a single cycle through every line, which no prefetcher can predict.
\omcode{code/firm/memkit/cpp/firm_memkit.hpp}{52}{72}{Building and following a pointer-chase ring.}\label{lst:ll:memory-hierarchy-and-caches:chase}
\item \textbf{\omterm{def:ll:memory-hierarchy-and-caches:tlb}{Huge pages} on request.} \texttt{Buffer} maps anonymous memory, asks for transparent \omterm{def:ll:memory-hierarchy-and-caches:tlb}{huge pages} with
\texttt{madvise}, and writes every page once so that no page fault happens later on the hot path (chapter~6).
\omcode{code/firm/memkit/cpp/firm_memkit.hpp}{27}{36}{A buffer with huge pages and pre-faulting.}\label{lst:ll:memory-hierarchy-and-caches:buffer}
\item \textbf{A counter per line.} \texttt{CachePadded} aligns its value on a line of its own; two threads increment
their counters with relaxed atomic additions, once padded and once packed in one structure.
\omcode{code/firm/memkit/cpp/firm_memkit.hpp}{22}{26}{A value alone on its cache line.}\label{lst:ll:memory-hierarchy-and-caches:padded}
\item \textbf{Run} \texttt{python bench\_mem.py} (about a minute): the staircase with 4~KiB and \omterm{def:ll:memory-hierarchy-and-caches:tlb}{huge pages} and in address
order, the huge-page comparison over 256~MiB, and the false-sharing test on CPUs 2 and 4.
\end{tutsteps}
\textbf{What to change next.} Use a stride of 4\,096 bytes instead of 64 in the chase and find the \omterm{def:ll:memory-hierarchy-and-caches:miss}{working set} at which
the first level fails, far earlier than 48~KiB; pin the two false-sharing threads to the same CPU and explain the result.
\end{tutorial}

\section{Build: memory placement tools}\label{bld:ll:memory-hierarchy-and-caches:memkit}

\begin{build}
\bfield{Purpose} The layout primitives every later component uses: counters and sequence numbers written by different
threads (rings, chapter~12; the feed handler's statistics, chapter~18) and buffers that must not fault or miss the TLB
on the hot path (books, journals).
\bfield{Interface} C++20 \texttt{firm::memkit}: \texttt{kLine}, \texttt{CachePadded<T>}, \texttt{Buffer(bytes,
huge)} with \texttt{data()}, \texttt{size()}, \texttt{huge\_requested()}, \texttt{chase\_ring(base, n, stride, random,
seed)}, \texttt{chase(start, steps)}. Rust \texttt{firm\_memkit}: \texttt{CachePadded<T>}, \texttt{chase\_ring},
\texttt{chase} on a byte buffer.
\bfield{Rules} Padded values occupy whole lines; buffers are page-aligned, rounded up to the page size requested, and
pre-faulted at construction; a failed huge-page request falls back to base pages and says so.
\bfield{Acceptance tests} \texttt{code/firm/memkit/}: sizes and alignments of padded values, neighbouring padded
values one line apart, a 2~MiB-rounded huge buffer, and the ring as a single cycle in both orders, in C++ and Rust.
\bfield{Stretch} Report the \omterm{def:ll:memory-hierarchy-and-caches:tlb}{huge pages} actually obtained, from \texttt{/proc/self/smaps}; allocate from reserved 1~GiB
pages when the host has them.
\end{build}

\begin{omsources}
\item U.~Drepper, ``What every programmer should know about memory'', LWN.net, 2007.
\item M.~Papamarcos and J.~Patel, ``A low-overhead coherence solution for multiprocessors with private cache memories'',
ISCA 1984.
\item Linux kernel documentation, ``Transparent hugepage support''.
\item Chips and Cheese, ``Intel's Redwood Cove: baby steps are still steps'' (2024) and ``Update on Meteor Lake DRAM
latency measurements'' (2024).
\end{omsources}

\section{Exercises}

\begin{exercise}[$\star$]\label{exo:ll:memory-hierarchy-and-caches:1}
In a cache with 64-byte lines and 64 sets, which set does address \texttt{0x7f3a12345e48} map to, and what is its offset
within the line?
\end{exercise}

\begin{exercise}[$\star$]\label{exo:ll:memory-hierarchy-and-caches:2}
A 2~MiB, 16-way cache has 64-byte lines. How many sets does it have, and at what stride do addresses alias to the same
set?
\end{exercise}

\begin{exercise}[$\star$]\label{exo:ll:memory-hierarchy-and-caches:3}
A process uses 600~MiB of hot data. How many TLB entries does it need with 4~KiB pages, and with 2~MiB pages?
\end{exercise}

\begin{exercise}[$\star\star$]\label{exo:ll:memory-hierarchy-and-caches:4}
A message structure holds an 8-byte sequence number written by the feed thread and a 4-byte counter written by the
strategy thread, next to each other. What happens, and how do you fix it at a cost of how many bytes?
\end{exercise}

\begin{exercise}[$\star\star$]\label{exo:ll:memory-hierarchy-and-caches:5}
From \cref{fig:ll:memory-hierarchy-and-caches:stairs}, how long does a book update that makes four dependent random
loads take when its data fit in the second level, and when they come from memory?
\end{exercise}

\begin{exercise}[$\star\star$]\label{exo:ll:memory-hierarchy-and-caches:6}
Why does the address-order chase stay fast far beyond the second level, while the random chase does not? What does that
say about linked lists and balanced trees on a hot path?
\end{exercise}

\begin{exercise}[$\star\star\star$]\label{exo:ll:memory-hierarchy-and-caches:7}
\emph{Coding.} Rewrite the Rust \texttt{chase} of \texttt{firm\_memkit} with raw pointers and \texttt{unsafe}, as the C++
does, and compare both with the index version on a 64~MiB ring. Is the bounds check visible in the time per step? Why or
why not?
\end{exercise}

\begin{exercise}[$\star\star\star$]\label{exo:ll:memory-hierarchy-and-caches:8}
\emph{Find the flaw.} ``We benchmarked our order-book update at \qty{20}{\nano\second} by applying the same add and
delete messages to one instrument ten million times in a loop.''
\end{exercise}

\section{Problem: The Staircase}

\begin{problem}[{Weekend problem --- how many books fit}]\label{pb:ll:memory-hierarchy-and-caches:1}
A market-making process quotes a set of instruments. For each it keeps a dense price ladder of 256 levels per side,
each level 32 bytes (price, quantity, order count, flags), plus 4~KiB of per-instrument state. The team wants every
instrument's hot data in the second-level cache of the core that runs the strategy.

\medskip\noindent\textbf{Part I --- The measurement.}
\begin{enumerate}
\item Read the latency of the first level, the second level and memory from
\cref{fig:ll:memory-hierarchy-and-caches:stairs}.
\item Where are the steps, and what capacities do they suggest?
\item Why does the third-level step end at about eight megabytes, not at 24~MiB?
\item What sets the height of the step between the second level and memory?
\end{enumerate}

\medskip\noindent\textbf{Part II --- The layout.}
\begin{enumerate}[resume]
\item How many bytes does one instrument's hot data occupy?
\item How many \omterm{def:ll:memory-hierarchy-and-caches:line}{cache lines}?
\item If three quarters of a 2~MiB second level is available to the books, how many instruments fit?
\item How many fit in a 48~KiB first level, at most?
\end{enumerate}

\medskip\noindent\textbf{Part III --- When they do not fit.}
\begin{enumerate}[resume]
\item The desk wants 300 instruments. What fraction of updates will miss the second level if instruments are updated
uniformly at random, and the second level holds the most recent ones?
\item If an update makes three dependent loads, what is its average memory time with that miss rate, taking the
second-level and memory latencies of Part~I?
\item What would \omterm{def:ll:memory-hierarchy-and-caches:tlb}{huge pages} change, and why especially here?
\item Name two layout changes that shrink the hot data.
\end{enumerate}

\medskip\noindent\textbf{Part IV --- The verdict.}
\begin{enumerate}[resume]
\item State the \emph{named result}: the latency of each level on this laptop and the number of instruments whose book
fits in the second level.
\item Why is 32 bytes a good size for a level, and 48 a bad one?
\item Should the ladder store levels as structures of four fields or four arrays of one field? When does each win?
\item What does the chase measure that a sequential benchmark hides?
\item Why must the benchmark visit every line once per lap?
\item What is the cost of one false-shared counter in the book, per update?
\item How would you confirm on the production host that the books stay in the second level?
\item In one sentence: what decides the speed of an order book?
\end{enumerate}
\end{problem}

\section{Interview questions}

\begin{interviewq}[$\star$ \iqroles{developer}]\label{iq:ll:memory-hierarchy-and-caches:1}
Why is iterating over a \texttt{std::vector} faster than over a \texttt{std::list} of the same elements?
\end{interviewq}

\begin{interviewq}[$\star\star$ \iqroles{developer}]\label{iq:ll:memory-hierarchy-and-caches:2}
What is \omterm{def:ll:memory-hierarchy-and-caches:false}{false sharing}? How would you detect it and fix it?
\end{interviewq}

\begin{interviewq}[$\star\star$ \iqroles{developer}]\label{iq:ll:memory-hierarchy-and-caches:3}
What is a TLB, and why do trading systems use \omterm{def:ll:memory-hierarchy-and-caches:tlb}{huge pages}?
\end{interviewq}

\begin{interviewq}[$\star\star$ \iqroles{developer}]\label{iq:ll:memory-hierarchy-and-caches:4}
How would you measure the latency of each level of a machine's cache hierarchy?
\end{interviewq}

\begin{interviewq}[$\star\star$ \iqroles{developer}]\label{iq:ll:memory-hierarchy-and-caches:5}
Two arrays of 4\,096 doubles are processed element by element together and the loop is unexpectedly slow. What might be
happening?
\end{interviewq}

\begin{interviewq}[$\star\star\star$ \iqroles{developer}]\label{iq:ll:memory-hierarchy-and-caches:6}
Design the memory layout of an order-id to order lookup for a feed with a million live orders, for the lowest p99.
\end{interviewq}
